\chapter{İstatistiksel Düşünme, Veri ve Araştırma Süreci}
\label{ch:research}
\label{ch:population-data}
\index{istatistiksel düşünme@İstatistiksel düşünme}

\begin{hedefler}
Bu bölümün sonunda okuyucu:
\begin{itemize}
  \item bir PDR problemini betimsel, karşılaştırmalı veya ilişkisel araştırma sorusuna dönüştürebilecek,
  \item yapı, değişken ve işlemsel tanım arasındaki ilişkiyi açıklayabilecek,
  \item evren, örneklem, analiz birimi, parametre ve istatistiği ayırt edebilecek,
  \item ölçme düzeyi, araştırma tasarımı ve örnekleme yöntemine uygun bir analiz planı kurabilecek,
  \item veri toplama ve raporlama kararlarını etik ve eleştirel biçimde değerlendirebilecektir.
\end{itemize}
\end{hedefler}

\begin{hazirlik}
``Öğrencilerin kaygısı yüksektir'' cümlesini ölçülebilir bir araştırma
sorusuna dönüştürün. Kimlerin, hangi kaygının ve hangi ölçümün kastedildiğini
yazın.
\end{hazirlik}

\section{PDR vakası: Üniversiteye uyum programı}

Bir üniversitenin psikolojik danışma birimi, birinci sınıf öğrencileri için
dört haftalık bir uyum programı hazırlamıştır. Birim yöneticisi programın
``öğrencilerin uyumunu artırıp artırmadığını'' bilmek istemektedir. İlk bakışta
basit görünen bu cümle çok sayıda araştırma kararı içerir:

\begin{itemize}
  \item Uyum hangi boyutları kapsar ve nasıl ölçülecektir?
  \item Sonuç hangi öğrenci grubuna genellenecektir?
  \item Programa katılanlarla katılmayanlar başlangıçta karşılaştırılabilir midir?
  \item Uyum ne zaman ölçülecektir: program öncesinde, sonrasında veya her ikisinde mi?
  \item Gözlenen değişim programa mı, zamana mı, başka deneyimlere mi bağlanabilir?
\end{itemize}

İstatistiksel düşünme, bu soruları bir yazılım menüsü seçmeden önce
yanıtlamaktır. Bölüm boyunca bu vaka, soyut kavramları somut bir araştırma
planına bağlamak için kullanılacaktır.

\section{İstatistik neden yalnızca hesaplama değildir?}
\index{istatistik}
\index{veri üretim süreci}

İstatistik, eldeki sayılara bir formül uygulamaktan önce verinin nasıl
üretildiğini sorgulayan bir düşünme biçimidir. Sağlam bir analiz şu zinciri
izler:
\begin{center}
\begin{tikzpicture}[
  alt={Sağlam bir analizin beş adımı: araştırma sorusu, ölçülebilir değişkenler, veri toplama tasarımı, analiz ve bağlamsal yorum.},
  x=0.20\linewidth,
  adim/.style={draw=PrismBlue,fill=PrismLight,rounded corners=3pt,
    text width=0.155\linewidth,minimum height=1.2cm,inner sep=4pt,
    align=center,font=\small\sffamily},
  ok/.style={->,thick,PrismBlue}
]
  \node[adim] (soru) at (0,0) {Araştırma\\sorusu};
  \node[adim] (degisken) at (1,0) {Ölçülebilir\\değişkenler};
  \node[adim] (tasarim) at (2,0) {Veri toplama\\tasarımı};
  \node[adim] (analiz) at (3,0) {Analiz};
  \node[adim] (yorum) at (4,0) {Bağlamsal\\yorum};
  \draw[ok] (soru.east)--(degisken.west);
  \draw[ok] (degisken.east)--(tasarim.west);
  \draw[ok] (tasarim.east)--(analiz.west);
  \draw[ok] (analiz.east)--(yorum.west);
\end{tikzpicture}
\end{center}
Zincirin ilk halkalarında yapılan bir hata, daha sonra kullanılan yöntem ne
kadar gelişmiş olursa olsun giderilemez. Örneğin yalnız gönüllü öğrencilerden
toplanan bir memnuniyet anketinde örnekleme yanlılığı varsa, çok küçük bir
$p$-değeri bu yanlılığı ortadan kaldırmaz \cite{longford2008,lohr2021}.

\subsection{Değişkenlik neden merkezîdir?}
\index{değişkenlik}

PDR verilerinde bireyler; yaşantı, sosyal destek, ölçüm hatası ve bağlamsal
koşullar nedeniyle birbirinden farklıdır. Aynı bireyin puanı bile farklı
zamanlarda değişebilir. İstatistik bu değişkenliği yok saymaz; kaynağını,
büyüklüğünü ve sonuç üzerindeki etkisini anlamaya çalışır.

Bir araştırmada gözlenen fark kabaca şu bileşenlerin birleşimi olarak
düşünülebilir:
\[
\text{gözlenen farklılık}
=\text{sistematik etki}+\text{bireysel farklılık}+\text{ölçüm ve örnekleme hatası}.
\]
İyi bir tasarım sistematik etkiyi diğer kaynaklardan ayırmayı kolaylaştırır;
istatistiksel yöntem ise kalan belirsizliği görünür kılar.

\begin{sezgibox}
İstatistik ``veride fark var mı?'' sorusundan önce ``bu fark hangi süreçlerle
oluşmuş olabilir?'' sorusunu sorar. Hesaplama, veri üretim sürecinin yerini
tutmaz.
\end{sezgibox}

\section{Araştırma probleminden araştırma sorusuna}
\index{araştırma problemi}
\index{araştırma sorusu}

``Ergenlerin yalnızlığı'', ``okul iklimi'' veya ``kariyer kararsızlığı'' birer
araştırma alanıdır; tek başına yanıtlanabilir araştırma sorusu değildir.
İyi bir soru hedef grubu, değişkenleri ve amaçlanan karşılaştırma ya da ilişkiyi
açık eder.

\begin{table}[htbp]
\centering
\caption{Araştırma amacına göre soru örnekleri.}
\begin{tabular}{@{}p{0.20\textwidth}p{0.66\textwidth}@{}}
\toprule
Soru türü & PDR örneği \\
\midrule
Betimsel & Birinci sınıf öğrencilerinin akademik kaygı puanları nasıl dağılmaktadır? \\
Gruplar arası & Uyum programına katılan ve katılmayan öğrencilerin son-test uyum puanları farklı mıdır? \\
İlişkisel & Algılanan sosyal destek ile yalnızlık puanı arasında ilişki var mıdır? \\
Yordayıcı & Haftalık uyku süresi akademik kaygı puanını ne ölçüde yordamaktadır? \\
\bottomrule
\end{tabular}
\end{table}

Betimsel soru örneklemdeki düzeyi veya dağılımı açıklar. Karşılaştırmalı soru
gruplar ya da zamanlar arasındaki farkı, ilişkisel soru değişkenlerin birlikte
değişimini, yordayıcı soru ise bir değişkenden diğerine ilişkin öngörüyü
inceler. İlişki ve yordama dili, tek başına nedensellik anlamına gelmez.

\subsection{Araştırma sorusu ve hipotez}
\index{hipotez}

Araştırma sorusu neyin öğrenilmek istendiğini belirtir. İstatistiksel hipotez
ise evren parametresi hakkında sınanabilir bir önerme kurar. Örneğin
``program gruplarının ortalama uyum puanları farklı mıdır?'' sorusu için
\[
H_0:\mu_{\text{program}}-\mu_{\text{karşılaştırma}}=0,
\qquad
H_1:\mu_{\text{program}}-\mu_{\text{karşılaştırma}}\ne0
\]
yazılabilir. Hipotezler veriyi gördükten sonra sonuca uygun biçimde
değiştirilmemelidir.

\section{Temel kavramlar}
\label{sec:population-parameter}
\index{evren}
\index{örneklem}
\index{parametre}
\index{istatistik}

\begin{description}
  \item[Evren (anakütle)] Araştırma sonucunun genellenmek istendiği bütün birimlerdir.
  \item[Örneklem] Evrenden gözlenen veya ölçülen birimler kümesidir.
  \item[Parametre] Evrene ait, çoğu zaman bilinmeyen sayısal özelliktir.
  \item[İstatistik] Örneklem verisinden hesaplanan sayısal özettir.
  \item[Değişken] Birimden birime farklı değer alabilen ölçülen özelliktir.
  \item[Analiz birimi] Veri tablosundaki her satırın temsil ettiği kişi, grup, kurum veya ölçümdür.
  \item[Gözlem] Bir analiz birimi için kaydedilen değişken değerlerinin bütünüdür.
\end{description}

Bir üniversitedeki bütün birinci sınıf öğrencilerinin matematik kaygısı
ortalaması bir parametre, seçilen 120 öğrencinin ortalaması ise bir
istatistiktir. İstatistik, parametre hakkında bilgi edinmek için kullanılır.

\begin{etkinlik}
Uyum programı vakasında hedef evreni, erişilebilir evreni, örneklemi, analiz
birimini, temel parametreyi ve bu parametreyi tahmin edecek istatistiği ayrı
ayrı yazın. Hedef evren ile erişilebilir evrenin neden farklı olabileceğini
tartışın.
\end{etkinlik}

\section{Yapıdan ölçülebilir değişkene}
\index{yapı}
\index{operasyonel tanım}
\index{ölçme}

Kaygı, öz-yeterlik, psikolojik dayanıklılık ve okul aidiyeti doğrudan
gözlenemeyen kuramsal yapılardır. Araştırmacı önce yapının hangi anlamda
kullanıldığını kavramsal olarak açıklar, sonra onu gözlenebilir bir değişkene
dönüştüren işlemsel tanımı kurar.

\begin{table}[htbp]
\centering
\caption{Kuramsal yapıdan işlemsel tanıma geçiş.}
\begin{tabular}{@{}p{0.20\textwidth}p{0.31\textwidth}p{0.34\textwidth}@{}}
\toprule
Yapı & Kavramsal anlam & Olası işlemsel tanım \\
\midrule
Akademik kaygı & Akademik değerlendirmelerle ilişkili endişe ve gerilim & Geçerliği desteklenmiş ölçeğin toplam puanı \\
Okula aidiyet & Öğrencinin okul topluluğuna kabul ve bağlılık algısı & Aidiyet ölçeği ortalama puanı \\
Yardım arama & Psikolojik destek kaynağına yönelme davranışı & Son altı ayda danışma birimine başvuru: evet/hayır \\
Devamlılık & Öğrencinin programa katılım düzeyi & Katıldığı oturum sayısı: 0--4 \\
\bottomrule
\end{tabular}
\end{table}

Aynı yapı farklı biçimlerde işlemselleştirilebilir. Kaygıyı tek maddelik bir
soruyla, klinik görüşmeyle veya çok maddeli ölçekle ölçmek aynı veri kalitesini
üretmez. İşlemsel tanım, hem istatistiksel yöntemi hem de sonucun hangi anlamda
genellenebileceğini sınırlar.

\section{Değişken türleri ve ölçme düzeyleri}
\label{sec:population-data-types}
\index{değişken}
\index{ölçme düzeyi}
\index{nominal ölçek}
\index{sıralı ölçek}
\index{aralıklı ölçek}
\index{oranlı ölçek}

\begin{table}[htbp]
\centering
\caption{Ölçme düzeyleri ve uygun temel özetler.}
\begin{tabular}{p{0.18\textwidth}p{0.31\textwidth}p{0.35\textwidth}}
\toprule
Düzey & Örnek & Uygun özetler \\
\midrule
Nominal & Program türü, cinsiyet kategorisi & Frekans, yüzde, mod \\
Sıralı & Memnuniyet düzeyi, Likert maddesi & Frekans, medyan, yüzdelikler \\
Aralık & Standartlaştırılmış test puanı & Ortalama, standart sapma, fark \\
Oran & Süre, yaş, doğru cevap sayısı & Tüm aritmetik özetler ve oranlar \\
\bottomrule
\end{tabular}
\end{table}

Sayısal kod verilmesi bir değişkeni otomatik olarak nicel yapmaz. ``Hiç
katılmıyorum'' ile ``tamamen katılıyorum'' seçeneklerine 1--5 kodlarının
verilmesi kategoriler arasında bir sıra oluşturur; fakat tek bir madde için
ardışık kategoriler arasındaki uzaklıkların kesinlikle eşit olduğu sonucunu
garanti etmez.

\subsection{Değişkenlerin araştırmadaki rolleri}

Ölçme düzeyi değişkenin nasıl kaydedildiğini, araştırmadaki rol ise analizde
nasıl kullanıldığını anlatır. Bir değişken aşağıdaki rollerden birini veya
birkaçını üstlenebilir:

\begin{description}
  \item[Yanıt değişkeni] Açıklanmak, karşılaştırılmak veya yordanmak istenen sonuçtur.
  \item[Açıklayıcı değişken] Yanıtla ilişkisi incelenen grup, özellik veya ölçümdür.
  \item[Karıştırıcı değişken] Hem açıklayıcı değişkenle hem yanıtla ilişkili olup gözlenen ilişkiyi çarpıtabilir.
  \item[Kimlik değişkeni] Kayıtları ayırır; genellikle analizde sayısal büyüklük olarak kullanılmaz.
\end{description}

Uyum programında son-test uyum puanı yanıt, program grubu açıklayıcı değişken
olabilir. Programa gönüllü katılım başlangıç motivasyonuyla ilişkiliyse
motivasyon olası bir karıştırıcıdır.

\begin{mistakebox}
Öğrenci numarası rakamlardan oluşsa da nicel bir değişken değildir. Bu
numaraların ortalamasını hesaplamak araştırma açısından anlam taşımaz.
\end{mistakebox}

\section{Araştırma tasarımını tanımak}
\index{araştırma tasarımı}
\index{deneysel tasarım}
\index{gözlemsel araştırma}

\begin{description}
  \item[Kesitsel tasarım] Değişkenler belirli bir zaman kesitinde ölçülür.
  \item[Boylamsal tasarım] Aynı birimler zaman içinde tekrar ölçülür.
  \item[Gözlemsel tasarım] Araştırmacı var olan özellikleri gözler, müdahale atamaz.
  \item[Deneysel tasarım] Araştırmacı müdahaleyi düzenler ve ideal olarak birimleri koşullara rastgele atar.
\end{description}

Gözlemsel bir çalışmada programa gönüllü katılanlarla katılmayanların
karşılaştırılması seçim yanlılığı taşıyabilir. Rastgele atama, başlangıçta
bilinen ve bilinmeyen özellikleri dengeleme olasılığını artırır. Ancak rastgele
atama tek başına ölçüm kalitesi, kayıp gözlem veya etik sorunları çözmez.

\begin{mistakebox}
``Ön-test ve son-test arasında anlamlı fark bulundu; program kesinlikle
etkilidir'' sonucu her tasarımda savunulamaz. Karşılaştırma grubu yoksa zaman,
olgunlaşma, başka yaşantılar ve ölçüme alışma gibi alternatif açıklamalar
elenemeyebilir.
\end{mistakebox}

\section{İstatistiksel araştırma döngüsü}
\index{araştırma döngüsü}

\begin{enumerate}
  \item Araştırma problemi ve hedef evren açıkça tanımlanır.
  \item Değişkenler işlemsel olarak tanımlanır ve ölçme düzeyleri belirlenir.
  \item Örnekleme ve veri toplama yöntemi seçilir.
  \item Veri temizlenir; eksik, olanaksız ve aykırı değerler incelenir.
  \item Önce grafiksel ve betimsel, sonra uygun çıkarımsal analiz yapılır.
  \item Belirsizlik, etki büyüklüğü, sınırlılıklar ve bağlam birlikte raporlanır.
\end{enumerate}

Bu döngü doğrusal görünse de uygulamada yinelemelidir. İlk grafikler veri
toplama hatasını gösterebilir; ölçme aracının yapısı araştırma sorusunu
daraltabilir; varsayım incelemesi analiz planını değiştirebilir. Yapılan her
değişiklik gerekçesiyle kaydedilmelidir.

\section{Veri matrisi ve kod kitabı}
\index{veri matrisi}
\index{kod kitabı}
\index{analiz birimi}

İyi tasarlanmış bir veri dosyasında çoğunlukla her satır bir analiz birimini,
her sütun bir değişkeni temsil eder. Değişken adı, açıklaması, ölçme düzeyi,
kodları, birimi, olası değer aralığı ve eksik değer gösterimi kod kitabında
belirtilir.

\begin{table}[htbp]
\centering
\caption{Uyum programı için örnek veri matrisi.}
\begin{tabular}{@{}rrrrr@{}}
\toprule
\texttt{id} & \texttt{grup} & \texttt{uyum\_on} & \texttt{uyum\_son} & \texttt{oturum} \\
\midrule
101 & 1 & 42 & 55 & 4 \\
102 & 1 & 51 & 57 & 3 \\
103 & 0 & 48 & 49 & 0 \\
104 & 0 & 54 & 52 & 0 \\
\bottomrule
\end{tabular}
\end{table}

Burada \texttt{grup} için 1 program, 0 karşılaştırma grubunu; \texttt{oturum}
ise katılım sayısını göstermektedir. Bu dört satır analiz için yeterli değildir;
tablo yalnız veri yapısını göstermektedir.

\begin{lstlisting}
import pandas as pd

veri = pd.DataFrame({
    "id": [101, 102, 103, 104],
    "grup": [1, 1, 0, 0],
    "uyum_on": [42, 51, 48, 54],
    "uyum_son": [55, 57, 49, 52],
    "oturum": [4, 3, 0, 0]
})

print(veri.info())
print(veri.isna().sum())
print(veri.describe())
\end{lstlisting}

Bu ilk inceleme analiz sonucu üretmez; değişken türlerini, eksikleri ve değer
aralıklarını denetler. Kimlik değişkeninin ortalaması yazılım tarafından
hesaplansa bile araştırma açısından yorumlanmaz.

\section{Örnekleme ve yanlılık}
\index{örnekleme}
\index{yanlılık}

Basit rastgele örneklemede her birimin bilinen bir seçilme şansı vardır.
Tabakalı örneklemede önemli alt gruplar ayrı ayrı temsil edilir. Küme
örneklemesinde doğal gruplar, örneğin sınıflar veya okullar seçilir. Kolayda
örnekleme hızlıdır; ancak hedef evreni temsil etme gücü çoğu zaman sınırlıdır.

\begin{researchbox}
Yanlılık örneklem hacmi büyütülerek otomatik olarak giderilemez. Yanlı bir
seçim mekanizmasından gelen 10.000 gözlem, iyi tasarlanmış daha küçük bir
örneklemden daha yanıltıcı olabilir.
\end{researchbox}

Örnekleme yöntemleri bir sonraki bölümde ayrıntılı ele alınacaktır. Bu aşamada
temel ilke, örneklemden evrene yapılacak her genellemenin seçim mekanizmasına
dayandığını fark etmektir.

\section{Etik ve eleştirel veri okuryazarlığı}
\index{araştırma etiği}
\index{veri okuryazarlığı}

İstatistiksel sonuçlar kişisel verilerin korunması, katılımcı onamı, eksik
verilerin açıklanması ve analiz kararlarının şeffaflığıyla birlikte
değerlendirilir. Yalnız anlamlı sonuçları seçerek raporlamak, korelasyonu
nedensellik gibi sunmak veya analiz sonrasında hipotezi önceden belirlenmiş
gibi göstermek bilimsel güvenilirliği zedeler.

PDR verileri ruh sağlığı, travma, aile ilişkileri veya yardım arama davranışı
gibi hassas bilgiler içerebilir. Etik bir veri yönetimi en az şu soruları
yanıtlamalıdır:

\begin{itemize}
  \item Katılımcı neye onay verdi ve çalışmadan çekilme hakkını biliyor mu?
  \item Kimliği belirleyebilecek değişkenler analiz dosyasından ayrıldı mı?
  \item Küçük alt grupların tablolarda yeniden tanınma riski var mı?
  \item Veriye kimler, hangi amaçla ve ne kadar süre erişebilecek?
  \item Bulgular damgalayıcı veya suçlayıcı bir dille sunuluyor mu?
\end{itemize}

\subsection{Şeffaflık ve yeniden üretilebilirlik}
\index{şeffaflık}
\index{yeniden üretilebilirlik}

Araştırma soruları, dışlama ölçütleri ve temel analiz planı mümkün olduğunca
veri görülmeden önce yazılmalıdır. Veri temizleme kararları, kullanılan yazılım
ve analiz adımları kaydedilmelidir. Yeniden üretilebilirlik katılımcı verisini
herkese açmak anlamına gelmez; hassas veriler korunurken anonimleştirilmiş
örnek veri, kod ve karar günlüğü paylaşılabilir.

\section{Çözümlü araştırma planı}

Uyum programı vakası için savunulabilir bir başlangıç planı şöyledir:

\begin{enumerate}
  \item \textbf{Soru:} Program ve karşılaştırma gruplarının uyum puanındaki değişimleri farklı mıdır?
  \item \textbf{Evren:} İlgili üniversitedeki yeni kayıtlı birinci sınıf öğrencileri.
  \item \textbf{Analiz birimi:} Programa uygun her öğrenci.
  \item \textbf{Yanıt:} Geçerliği desteklenmiş uyum ölçeğinin son-test veya değişim puanı.
  \item \textbf{Açıklayıcı değişken:} Program grubu; 1 program, 0 karşılaştırma.
  \item \textbf{Tasarım:} Ön-test--son-test ve karşılaştırma gruplu tasarım; atama yönteminin açıkça belirtilmesi gerekir.
  \item \textbf{İlk analiz:} Eksik veri ve değer aralığı kontrolü; gruplara göre grafik, ortalama, standart sapma, medyan ve IQR.
  \item \textbf{Çıkarım:} Tasarıma ve varsayımlara uygun grup/zaman karşılaştırması; etki büyüklüğü ve güven aralığı.
  \item \textbf{Sınırlılık:} Gönüllü katılım varsa seçilim etkisi ve nedensel yorum sınırı.
  \item \textbf{Etik:} Bilgilendirilmiş onam, gizlilik, destek gereksinimi belirlenen öğrenci için yönlendirme protokolü.
\end{enumerate}

\begin{outputguidebox}
Bir yazılım çıktısı araştırma planındaki boşlukları tamamlayamaz. Çıktıda
``anlamlı'' bir değer bulunması hedef evreni, ölçme aracının geçerliğini,
karıştırıcı değişkenleri veya etik raporlama sorumluluğunu otomatik olarak
çözmez.
\end{outputguidebox}

\section{Literatür veri setiyle IBM SPSS laboratuvarı}
\index{Student Performance veri seti}
\index{SPSS!veri içe aktarma}
\index{veri sözlüğü}

Bu ve izleyen bölümlerde yöntemler arasındaki bağı görünür kılmak için Cortez
ve Silva'nın ortaöğretimde öğrenci başarısını incelediği çalışmayla ilişkili
\emph{Student Performance} veri setinin matematik dersi dosyası
kullanılacaktır \cite{cortezsilva2008,cortez2008data}. Veri, iki Portekiz
okulundaki okul kayıtları ve öğrenci anketlerinden oluşturulmuştur. Matematik
dosyasında 395 öğrenci ve 33 değişken bulunur. Veri seti CC BY 4.0 lisansıyla
yayımlanmıştır; bu nedenle kaynak gösterilerek öğretim amacıyla kullanılabilir.

\begin{softwarebox}
UCI deposundan \texttt{student.zip} dosyasını indirip içindeki
\texttt{student-mat.csv} dosyasını çalışma klasörüne çıkarın. Kitapta veri
dosyasının tamamı yeniden basılmamış; lisans, özgün değişken adları ve kaynak
künyesi korunmuştur. Bu yaklaşım hem sürüm denetimini hem de atfı açık tutar.
\end{softwarebox}

\subsection{Bu bölümdeki araştırma görevi}

Bu ilk uygulamanın amacı bir hipotez sınamak değil, veri üretim süreci ile
değişkenlerin analizdeki rollerini tanımaktır. Öğrenci şu soruları
yanıtlamalıdır:

\begin{enumerate}
  \item Her satır ve her sütun neyi temsil etmektedir?
  \item Sayısal kodlanan değişkenlerin hangileri nominal, sıralı veya niceldir?
  \item Dönem notları \texttt{G1}, \texttt{G2} ve \texttt{G3} aynı öğrenciye ait olduğundan hangi bağımlılık yapısını taşır?
  \item Değer aralıkları ve eksik gözlemler veri sözlüğüyle uyumlu mudur?
  \item Gözlemsel veri hangi yorumlara izin verir, hangilerine vermez?
\end{enumerate}

\subsection{SPSS syntax ile içe aktarma ve denetim}

Aşağıdaki syntax noktalı virgülle ayrılmış CSV dosyasını okur, değişkenlerin
ölçme düzeylerini tanımlar ve ilk kalite kontrol tablolarını üretir. Dosya yolu
kullanıcının klasör yapısına göre düzenlenmelidir.

\begin{lstlisting}[language={}]
GET DATA
 /TYPE=TXT
 /FILE='student-mat.csv'
 /ENCODING='UTF8'
 /DELCASE=LINE
 /DELIMITERS=";"
 /QUALIFIER='"'
 /ARRANGEMENT=DELIMITED
 /FIRSTCASE=2
 /VARIABLES=
 school A2 sex A1 age F2.0 address A1 famsize A3 Pstatus A1
 Medu F1.0 Fedu F1.0 Mjob A8 Fjob A8 reason A10 guardian A6
 traveltime F1.0 studytime F1.0 failures F1.0
 schoolsup A3 famsup A3 paid A3 activities A3 nursery A3
 higher A3 internet A3 romantic A3
 famrel F1.0 freetime F1.0 goout F1.0 Dalc F1.0 Walc F1.0
 health F1.0 absences F3.0 G1 F2.0 G2 F2.0 G3 F2.0.

DATASET NAME StudentMath WINDOW=FRONT.

VALUE LABELS studytime
 1 '2 saatten az' 2 '2--5 saat'
 3 '5--10 saat' 4 '10 saatten fazla'.

VARIABLE LEVEL
 school sex address famsize Pstatus Mjob Fjob reason guardian
 schoolsup famsup paid activities nursery higher internet
 romantic (NOMINAL)
 Medu Fedu traveltime studytime failures famrel freetime
 goout Dalc Walc health (ORDINAL)
 age absences G1 G2 G3 (SCALE).

DISPLAY DICTIONARY.
FREQUENCIES VARIABLES=school sex studytime higher.
DESCRIPTIVES VARIABLES=age absences G1 G2 G3
 /STATISTICS=MEAN STDDEV MIN MAX.
\end{lstlisting}

Menüyle çalışılacaksa \textbf{File $>$ Import Data $>$ CSV Data} yolunda
ayraç olarak noktalı virgül seçilir. İçe aktarmadan sonra \textbf{Variable
View} ekranında dize genişlikleri, etiketler ve ölçme düzeyleri syntax ile
karşılaştırılır. Menü ayarları \textbf{Paste} düğmesiyle syntax'a
dönüştürülerek analiz kaydı saklanmalıdır.

\subsection{Çıktının erişilebilir özeti}

SPSS Viewer görüntüsü yerine, ekran okuyucuyla da izlenebilen yeniden kurulmuş
çıktı tablosu verilmiştir. Aşağıdaki değerler özgün
\texttt{student-mat.csv} dosyasındaki \textbf{Descriptives} sonucudur.

\begin{table}[htbp]
\centering
\caption{Student Performance matematik verisi için temel SPSS betimsel çıktısı.}
\begin{tabular}{@{}lrrrrr@{}}
\toprule
Değişken & Geçerli $n$ & Ortalama & SS & En küçük & En büyük \\
\midrule
Yaş & 395 & 16,70 & 1,28 & 15 & 22 \\
Devamsızlık & 395 & 5,71 & 8,00 & 0 & 75 \\
Birinci dönem notu (G1) & 395 & 10,91 & 3,32 & 3 & 19 \\
İkinci dönem notu (G2) & 395 & 10,71 & 3,76 & 0 & 19 \\
Yıl sonu notu (G3) & 395 & 10,42 & 4,58 & 0 & 20 \\
\bottomrule
\end{tabular}
\end{table}

\subsection{Çıktının analizi}

Bütün değişkenlerde geçerli gözlem sayısının 395 olması, UCI veri sözlüğündeki
``eksik değer yok'' bilgisiyle uyumludur. Bu sonuç ölçüm hatası veya yanlılık
bulunmadığını göstermez; yalnız dosyada sistem eksik değer olmadığını gösterir.
Notların gözlenen aralıkları tanımlanan 0--20 sınırları içindedir. G1'in en
küçük değerinin 3, G2 ve G3'ün en küçük değerlerinin 0 olması veri hatası
kanıtı değildir; dönemler farklı ölçüm zamanlarıdır.

\texttt{studytime} değişkeni 1--4 sayılarıyla kodlanmış olsa da süre
kategorileri arasındaki aralıklar eşit değildir; bu nedenle ölçme düzeyi
\emph{sıralı} olarak tanımlanmıştır. Buna karşılık G3, bu kitapta nicel yanıt
değişkeni olarak ele alınacaktır. G1, G2 ve G3 aynı öğrenciye ait tekrarlı
notlardır; 1.185 bağımsız öğrenci gözlemi gibi yorumlanamaz.

Son olarak bu veri gözlemseldir. Örneğin çalışma süresi ile başarı arasında
ilişki bulunması, daha uzun çalışmanın tek başına daha yüksek nota neden
olduğunu kanıtlamaz. Önceki başarı, okul, aile özellikleri ve başka değişkenler
gözlenen ilişkiye katkıda bulunabilir. Sonraki bölümlerde üretilecek her SPSS
çıktısı bu tasarım sınırı içinde yorumlanacaktır.

\section{R ile çözümlü uygulama}
\label{sec:r-b01}

\textbf{Soru.} Bölümdeki beş öğrencilik öğretim örneğini R'de oluşturun.
Gözlem birimini, değişken türlerini, program frekanslarını ve ortalama başarıyı
belirleyin. Bu küçük veri gerçek bir araştırmadan alınmış değildir.

\begin{lstlisting}[style=prismR]
veri <- data.frame(
  devam_saati = c(8, 10, 7, 12, 9),
  basari = c(68, 75, 64, 83, 72),
  program = factor(c("A", "A", "B", "B", "A"))
)
str(veri)
summary(veri)
table(veri$program)
prop.table(table(veri$program))
colSums(is.na(veri))
c(ogrenci = nrow(veri),
  ortalama_devam = mean(veri$devam_saati),
  ortalama_basari = mean(veri$basari))
\end{lstlisting}

\begin{outputguidebox}
\textbf{Çözüm ve kontrol değerleri.} Her satır bir öğrencidir; veri
5 satır ve 3 değişken içerir. \texttt{devam\_saati} ve \texttt{basari}
nicel, \texttt{program} nominal kategoriktir. A programında 3, B programında
2 öğrenci vardır. Ortalama devam 9,2 saat, ortalama başarı 72,4 puandır;
hiçbir sütunda eksik değer yoktur.
\end{outputguidebox}

\textbf{Yorum.} \texttt{factor} program kodlarının kategori olarak
saklanmasını sağlar. Bir program etiketinin ortalamasını almak anlamlı
değildir. Beş satırın özetlenmesi, devamın başarıya neden olduğunu veya
öğrencilerin bir evreni temsil ettiğini göstermez.

\textbf{Pekiştirme ve yanıt.} Programların örneklemdeki paylarını
\texttt{prop.table} çıktısından okuyun: A için 0,60, B için 0,40 elde edilir. Bunlar evren oranı değil örneklem oranlarıdır.
\section{Bölüm sonu çözümlü problemler}

\begin{enumerate}
  \item \textbf{Araştırma sorusunu yapılandırma.} Bir danışma merkezi, ``akran desteğinin öğrencilerin okula uyumunu artırıp artırmadığını'' incelemek istiyor. Hedef evreni, yanıtı, açıklayıcı değişkeni ve uygun temel tasarımı belirtin.

  \textbf{Çözüm:} Hedef evren ilgili kurumun programa uygun öğrencileridir. Yanıt, geçerliği desteklenmiş bir uyum ölçeğinin son-test puanı; açıklayıcı değişken program grubudur. En güçlü temel tasarım, başlangıç ölçümü bulunan ve öğrencilerin program ile karşılaştırma koşullarına rastgele atandığı ön-test--son-test tasarımıdır. Rastgele atama yapılamıyorsa sonuç ilişki diliyle sunulmalı ve başlangıç farklılıkları incelenmelidir.

  \item \textbf{Parametre ve istatistik.} Üniversitedeki 4.000 birinci sınıf öğrencisinin yalnızlık ortalaması araştırılmaktadır. Rastgele seçilen 160 öğrencinin ortalaması 42,8 bulunmuştur. Parametre ve istatistiği ayırın.

  \textbf{Çözüm:} Parametre, 4.000 öğrencinin bilinmeyen evren ortalaması $\mu$'dür. Örneklemden hesaplanan $\bar{x}=42{,}8$ ise bu parametrenin istatistiksel tahminidir. 160 sayısı örneklem hacmi, 4.000 sayısı evren büyüklüğüdür.

  \item \textbf{Nedensel yorum sınırı.} Programa gönüllü katılan öğrencilerin kaygı ortalaması katılmayanlardan düşük bulunmuştur. ``Program kaygıyı azalttı'' sonucu neden doğrudan kurulamaz?

  \textbf{Çözüm:} Gönüllü katılım seçim yanlılığı yaratabilir. Motivasyon, başlangıç kaygısı, zaman uygunluğu veya yardım arama eğilimi hem katılımla hem sonuçla ilişkili olabilir. Zaman sırası, uygun karşılaştırma, atama süreci ve kayıplar değerlendirilmeden gözlenen fark nedensel etki olarak yorumlanamaz.
\end{enumerate}

\bolumkaynaklari{\cite{degroot2012,longford2008,lohr2021,anscombe1973,cortezsilva2008,cortez2008data}}

\section*{Hatırla--Uygula--Yansıt}

\begin{enumerate}
\renewcommand{\labelenumi}{\textbf{\Alph{enumi}.}}
  \item \textbf{Hatırla:} Evren, örneklem, parametre, istatistik ve analiz birimi kavramlarını tek bir PDR örneği üzerinde açıklayın.
  \item \textbf{Dönüştür:} ``Lise öğrencilerinde kariyer kararsızlığı'' konusunu betimsel, karşılaştırmalı ve ilişkisel üç soruya dönüştürün.
  \item \textbf{İşlemselleştir:} Okul aidiyeti yapısı için iki farklı işlemsel tanım önerin; güçlü ve sınırlı yönlerini karşılaştırın.
  \item \textbf{Sınıflandır:} Sınıf düzeyi, yaş, ölçek maddesi, toplam ölçek puanı ve danışmanlığa başvuru değişkenlerinin ölçme düzeylerini gerekçelendirin.
  \item \textbf{Eleştir:} Yalnız gönüllü öğrencilerle yürütülen çevrim içi kaygı anketinin hedef evrene genellenmesindeki sorunları yazın.
  \item \textbf{Tasarla:} Bir akran zorbalığı önleme programı için soru, evren, değişken, tasarım, örnekleme ve analiz birimlerini içeren kısa plan hazırlayın.
  \item \textbf{Etik karar:} Çok küçük bir alt grupta yüksek riskli davranış oranını raporlamanın gizlilik ve damgalama risklerini tartışın.
  \item \textbf{Yansıt:} Güçlü bir istatistiksel yöntemin zayıf veri üretim sürecini neden telafi edemeyeceğini açıklayın.
\end{enumerate}